\section{Exact binary profiles beyond the previous threshold}\label{sec:binary-profile}
For $t\ge1$ and $r>1$, write $B_2(t,r)$ for the maximum number of
geometric roots in $v_2(X-1)\ge r$ of a nonzero polynomial with at most
$t$ nonzero terms. Coefficients may lie in $\mathbb C_2$; allowing only
rational coefficients gives the same maxima below. Roots are counted with
multiplicity for the upper bounds, and all lower bounds have distinct roots.
Write $B_2^\circ(t)$ for the corresponding maximum in $v_2(X-1)>1$.

\begin{theorem}[Five, six, and seven terms over the binary disk]\label{thm:binary-profile}
For every real $r>1$,
\[
 B_2(5,r)=4,
 \qquad
 B_2(6,r)=
 \begin{cases}8,&1<r\le4/3,\\5,&r>4/3,\end{cases}
 \qquad
 B_2(7,r)=
 \begin{cases}8,&1<r\le3/2,\\6,&r>3/2.\end{cases}
\]
Consequently
\[
 B_2^\circ(5)=4,\qquad B_2^\circ(6)=B_2^\circ(7)=8.
\]
Each displayed maximum is attained by an integer polynomial with exactly
$t$ terms and with all of its roots in the indicated disk simple.
\end{theorem}

\begin{proof}
The cutoff theorem has $D=2$, $H=1$, $M=1$, $d=3$ and $N=8$ for
$t=5,6,7$. After normalizing the first $t$ moments to be integral with
at least one unit, any minimizing index $k\ge t$ has
$s_2(k)\le1$ and $v_2(M_k)<1$. Thus the only such index is $k=8$.
The coefficients are weighted by
\[
 w_r(k)=rk-v_2(k!)=(r-1)k+s_2(k).
\]
For $t=6$, the largest of the early weights is $w_r(5)=5r-3$ for
$r>1$. Since $w_r(8)=8r-7$, the index $8$ is impossible if $r>4/3$.
For $t=7$, the largest early weight is $w_r(6)=6r-4$, so the index $8$
is impossible if $r>3/2$. This proves the upper bounds of $8$ and the
stronger bounds beyond those respective endpoints.

For $t=5$, we exclude the index $8$ at every $r>1$. In the exponent
algebra modulo two, write
\[
 \overline A=\mathbb F_2[T]/\bigl(T^{a}(T-1)^{5-a}\bigr).
\]
If $1\le a\le4$, then $T^8=T^4$ in $\overline A$: on either primary
factor the fourth and eighth powers are the corresponding scalar zero
or one. Hence $M_8\equiv M_4\pmod2$. Because a minimizing $M_8$ would
have valuation below one, the two moments would have equal valuation.
But $w_r(8)-w_r(4)=4(r-1)>0$, a contradiction.
If $a=5$, then $T^8=0$, making $v_2(M_8)\ge1$, again impossible.
If $a=0$, then $T^8=1$, and comparison with $M_0$ gives the same
contradiction because $w_r(8)>0$.
This proves the upper bound four. Allowing fewer than five terms cannot
increase it: one may pad a support by unused columns before applying
the argument. The same padding convention covers the earlier bounds.

For sharpness in the two branches with value eight, use
\begin{align*}
 F_6(X)&=5X^8-28X^6+70X^4-140X^2+128X+65501,\\
 F_7(X)&=3X^8-28X^6+210X^4-448X^3+420X^2-192X+65571.
\end{align*}
Their exact expansions at one are
\begin{align*}
 F_6(1+Y)&=2^{16}+112Y^5+112Y^6+40Y^7+5Y^8,\\
 F_7(1+Y)&=2^{16}+56Y^6+24Y^7+3Y^8.
\end{align*}
The lower Newton polygon of the first has vertices $(0,16),(5,4),(8,0)$.
It gives five roots of valuation $12/5$ and three of valuation $4/3$.
The second has vertices $(0,16),(6,3),(8,0)$, giving six roots of
valuation $13/6$ and two of valuation $3/2$.
Thus both endpoint inequalities retain all eight roots.
These roots are simple. Indeed the following explicit identities hold:
\[
 \begin{split}
 &(X^6+X^3-1)F_6+
 (X^7+X^5+X^4+X^3+X^2+X)F_6'\equiv1\pmod3,\\
 &X^5F_7+(-2X^6-2X^4+X^2+2)F_7'\equiv1\pmod5.
 \end{split}
\]
In each identity the leading coefficient of $F_i$ is a unit modulo the
indicated prime, so coprimality after reduction implies squarefreeness
over $\mathbb Q$.

For the remaining branches, take $(X-1)^{t-1}-2^K$ with $K$ a sufficiently
large positive integer, specifically $K/(t-1)\ge r$. This polynomial has
exactly $t$ nonzero terms and $t-1$ distinct roots, each of depth
$K/(t-1)$. For the open disk choose a strict inequality. All assertions
for real $r$ also follow directly from these witnesses and from the
weight comparisons, or by rational approximation in the cutoff theorem.
\end{proof}

The five-term value stays at four, while the sixth permitted term doubles
the open-disk maximum to eight. The seventh term then increases the
closed depth through which eight roots can be retained, from $4/3$ to
$3/2$, without increasing the open-disk maximum. These are geometric
root counts in $\mathbb C_2$, not counts of roots lying in $\mathbb Q_2$.
The earlier repository result through $t=p+2$ only reaches four terms
at $p=2$; the present five-, six-, and seven-term formulas lie beyond
that range.
